\documentclass[10pt]{article}

\usepackage[margin=0.75in]{geometry}
\usepackage{amsmath,amsthm,amssymb,bm}
\usepackage{xcolor}
\usepackage{cancel}
\usepackage{graphicx}
\usepackage{changepage}
\usepackage{circuitikz}
\usepackage{pgfplots}
\usepackage{minted}
\usepackage{physics}
\usepackage{tikz}
\usepackage{tikz-3dplot}
\usepackage{hyperref}
\usepackage{siunitx}
\usepackage{fontspec}
\usepackage{relsize}
\usepackage{subfig}
\usepackage{todonotes}
\usepackage{contour}
\usepackage{multicol, multirow, booktabs}
\usepackage[breakable]{tcolorbox}
\usepackage[inline]{enumitem}

\theoremstyle{definition}
\newtheorem{problem}{Problem}
\newtheorem{soln}{Solution}

\pgfplotsset{compat=newest}
\usetikzlibrary{lindenmayersystems}
\usetikzlibrary{arrows}
\usetikzlibrary{calc}
\usetikzlibrary{positioning, fit}
\usetikzlibrary{3d, perspective}
\usetikzlibrary{math}
\usetikzlibrary{fadings}
\usetikzlibrary{angles,quotes} 
\usetikzlibrary{decorations.markings}
\usetikzlibrary{tikzmark}

\definecolor{incolor}{HTML}{303F9F}
\definecolor{outcolor}{HTML}{D84315}
\definecolor{cellborder}{HTML}{CFCFCF}
\definecolor{cellbackground}{HTML}{F7F7F7}
\newcommand{\ui}{\hat{i}}
\newcommand{\uj}{\hat{j}}
\newcommand{\uk}{\hat{k}}
\newcommand{\ux}{\hat{\mathbf{x}}}
\newcommand{\uy}{\hat{\mathbf{y}}}
\newcommand{\uz}{\hat{\mathbf{z}}}
\newcommand{\uv}[1]{\hat{\bv{{#1}}}}
\newcommand{\pr}[1]{{#1^\prime}}
\newcommand{\pd}[2][]{{\frac{\partial^{#1}}{\partial {{#2}^{#1}}}}}
\newcommand{\dif}[2][]{{\frac{d^{#1}}{d{{#2}^{#1}}}}}
\newcommand{\grd}{\nabla}

\pgfdeclarelayer{background}  
\pgfsetlayers{background,main}
\AtBeginDocument{\RenewCommandCopy\qty\SI}

\makeatletter
\newcommand{\boxspacing}{\kern\kvtcb@left@rule\kern\kvtcb@boxsep}
\makeatother
\newcommand{\prompt}[4]{
    \ttfamily\llap{{\color{#2}[#3]:\hspace{3pt}#4}}\vspace{-\baselineskip}
}

\newcommand{\thevenin}[2]{
  \begin{center}
    \begin{circuitikz} \draw
      (0,0) -- (2,0) to[battery1, l_=$V_{Th}\eq#1$] (2,2) 
      to[resistor, l_=$R_{Th}\eq#2$] (0,2)
      ;
      \draw [o-] (-.07,2.079);
      \draw [o-] (-.07,0.079);
    \end{circuitikz}
  \end{center}
}

\newcommand{\norton}[2]{
  \begin{center}
    \begin{circuitikz} \draw
      (0,0) -- (3,0) to[american current source, l_=$I_{N}\eq#1$] (3,2) -- (0,2) (2,0)
      to[resistor, l=$R_{N}\eq#2$] (2,2)
      ;
      \draw [o-] (-.07,2.079);
      \draw [o-] (-.07,0.079);
    \end{circuitikz}
  \end{center}
}

\DeclareMathOperator{\Div}{div}
\DeclareMathOperator{\Curl}{curl}
\DeclareMathOperator{\sgn}{sgn}
\DeclareMathOperator{\erfc}{erfc}

\newcommand{\highlight}[1]{\colorbox{yellow}{$\displaystyle #1$}}
\newcommand{\ti}[1]{\widetilde{#1}}
\renewcommand{\div}[1]{\bv{\nabla}\cdot #1}
\renewcommand{\curl}[1]{\bv{\nabla}\times #1}

\newfontface{\Kaufmann}{Kaufmann}
\DeclareTextFontCommand{\kf}{\Kaufmann}
\newcommand{\scriptr}{\fontsize{12pt}{12pt}\kf{r}\fontsize{10pt}{10pt}}

\newfontface{\KaufmannB}{Kaufmann Bold}
\DeclareTextFontCommand{\kfb}{\KaufmannB}
\newcommand{\bscriptr}{\fontsize{12pt}{12pt}\kfb{r}\fontsize{10pt}{10pt}}
\newcommand{\justif}[2]{&{#1}&\text{#2}}
\newcommand{\bv}[1]{\bm{\mathrm{#1}}}

\title{Physics 4700H: Assignment VII-A}
\author{Jeremy Favro (0805980) \\ Trent University, Peterborough, ON, Canada}
\date{\today}

\begin{document}
\maketitle

% PROBLEM 1
\begin{problem} Consider a system consisting of two particles.
\begin{enumerate}[label=(\alph*)]
  \item First, assume the particles are distinguishable, and label them $A$ and $B$. Suppose each
        particle can be in one of three states, where for $A$ these have energies $\alpha_1$, $\alpha_2$, and $\alpha_3$ whereas
        for $B$ they have energies $\beta_1$, $\beta_2$, and $\beta_3$. Assume that the state of $A$ is independent of
        the state of $B$ and vice versa. Show that the partition function for the combined system
        $Z_{tot} = Z_AZ_B$.
  \item Now, assume the particles are indistinguishable. Again, each particle can be in one of
        three states, with energies $\epsilon_1$, $\epsilon_2$, and $\epsilon_3$. The particles are not interacting, and there is no
        restriction on the number of particles in a given state. What is the partition function of the
        combined system in this case?
\end{enumerate}
\end{problem}
\begin{soln}~
  \begin{enumerate}[label=(\alph*)]
    \item First note that by our usual definition the individual partition functions are
          $$Z_A=e^{-\alpha_1/k_BT}+e^{-\alpha_2/k_BT}+e^{-\alpha_3/k_BT}$$
          and
          $$Z_B=e^{-\beta_1/k_BT}+e^{-\beta_2/k_BT}+e^{-\beta_3/k_BT}.$$
          The combined system of these particles is allowed to have energies corresponding to any sum of the individual particle energies.
          Defining the partition function as usual then this looks like
          $$Z_{tot}=\sum_i\sum_je^{-(\beta_i+\alpha_j)/k_BT}.$$
          We can write out all of these terms, there are $3\cdot \binom{3}{1}=9$ terms, three groups where we choose one $\beta$ and then select all $\alpha$s,
          \begin{align*}
            Z_{tot} & =e^{-(\beta_1+\alpha_1)/k_BT}+e^{-(\beta_1+\alpha_2)/k_BT}+e^{-(\beta_1+\alpha_3)/k_BT}       \\
                    & \quad+e^{-(\beta_2+\alpha_1)/k_BT}+e^{-(\beta_2+\alpha_2)/k_BT}+e^{-(\beta_2+\alpha_3)/k_BT}  \\
                    & \quad+e^{-(\beta_3+\alpha_1)/k_BT}+e^{-(\beta_3+\alpha_2)/k_BT}+e^{-(\beta_3+\alpha_3)/k_BT}.
          \end{align*}
          Now we look at $Z_AZ_B$ (I switched them at the first equals sign to make it look more like the above result),
          \begin{align*}
            Z_AZ_B & =(e^{-\beta_1/k_BT}+e^{-\beta_2/k_BT}+e^{-\beta_3/k_BT})(e^{-\alpha_1/k_BT}+e^{-\alpha_2/k_BT}+e^{-\alpha_3/k_BT}) \\
                   & =e^{-\beta_1/k_BT}(e^{-\alpha_1/k_BT}+e^{-\alpha_2/k_BT}+e^{-\alpha_3/k_BT})                                       \\
                   & \quad+e^{-\beta_2/k_BT}(e^{-\alpha_1/k_BT}+e^{-\alpha_2/k_BT}+e^{-\alpha_3/k_BT})                                  \\
                   & \quad+e^{-\beta_3/k_BT}(e^{-\alpha_1/k_BT}+e^{-\alpha_2/k_BT}+e^{-\alpha_3/k_BT})                                  \\
                   & =e^{-(\beta_1+\alpha_1)/k_BT}+e^{-(\beta_1+\alpha_2)/k_BT}+e^{-(\beta_1+\alpha_3)/k_BT}                            \\
                   & \quad+e^{-(\beta_2+\alpha_1)/k_BT}+e^{-(\beta_2+\alpha_2)/k_BT}+e^{-(\beta_2+\alpha_3)/k_BT}                       \\
                   & \quad+e^{-(\beta_3+\alpha_1)/k_BT}+e^{-(\beta_3+\alpha_2)/k_BT}+e^{-(\beta_3+\alpha_3)/k_BT},
          \end{align*}
          indeed, the same answer, so $Z_{tot}=Z_AZ_B$.
    \item In this case the partition functions for the two particles are identical,
          $$Z_1=e^{-\epsilon_1/k_BT}+e^{-\epsilon_2/k_BT}+e^{-\epsilon_3/k_BT}.$$
          The entire system then, as before, can have energies corresponding to any sum of $\epsilon$s
          and so has partition function
          $$Z_{tot}=\sum_{i}\sum_je^{-(\epsilon_i+\epsilon_j)/k_BT}.$$
          Like before, this will have nine terms. However, since $\epsilon_i$ could be the same as $\epsilon_j$, not all of the terms are unique.
          Writing them all out we see that
          \begin{align*}
            Z_{tot} & =e^{-(\epsilon_1+\epsilon_1)/k_BT}+e^{-(\epsilon_1+\epsilon_2)/k_BT}+e^{-(\epsilon_1+\epsilon_3)/k_BT}         \\
                    & \quad+e^{-(\epsilon_2+\epsilon_1)/k_BT}+e^{-(\epsilon_2+\epsilon_2)/k_BT}+e^{-(\epsilon_2+\epsilon_3)/k_BT}    \\
                    & \quad+e^{-(\epsilon_3+\epsilon_1)/k_BT}+e^{-(\epsilon_3+\epsilon_2)/k_BT}+e^{-(\epsilon_3+\epsilon_3)/k_BT}    \\
                    & =e^{-2\epsilon_1/k_BT}+e^{-2\epsilon_2/k_BT}+e^{-2\epsilon_3/k_BT}                                             \\
                    & \quad+2e^{-(\epsilon_1+\epsilon_2)/k_BT}+2e^{-(\epsilon_1+\epsilon_3)/k_BT}+2e^{-(\epsilon_2+\epsilon_3)/k_BT}
          \end{align*}
  \end{enumerate}
\end{soln}

% PROBLEM 2
\begin{problem}
Consider a classical (i.e. high temperature and/or low density) gas of $N$ indistinguishable
non-interacting spin-1 particles, each of mass $m$, confined to a two-dimensional box of size
$L^2$. What is the entropy at temperature $T$? Note that spin-1 particles have three possible
spin states, which can be labeled $\uparrow$, $0$, and $\downarrow$. There is no applied magnetic field in this
problem, so these spin states all have the same energy. You may define that spin energy to
be zero.
\end{problem}
\begin{soln}

\end{soln}
\end{document}